% Chapter 2 of Volume III, "Routes, filters and the target in six forms": additions to ch/nc_route.tex. Sources: misc/reviews/nick-volume3/vol3/chC_path_detail.tex 11-34, 36-62, 64-78, 80-98, 123-133 and ch9_path_to_rh.tex 11-29, 44-52; reports U12 (M6, M7, M9, M10, M12), verification W3 (V15). 17 Nov 2026.
\ifdefined\computation\else\theoremstyle{sc}\newtheorem{computation}[theorem]{Computation}\theoremstyle{scr}\fi
\ifdefined\openproblem\else\theoremstyle{sc}\newtheorem{openproblem}[theorem]{Open problem}\theoremstyle{scr}\fi
\ifdefined\conjecture\else\theoremstyle{sc}\newtheorem{conjecture}[theorem]{Conjecture}\theoremstyle{scr}\fi
% The text follows ch/nc_route.tex inside one chapter and extends its Sections route:sec:target and route:sec:filters by reference. The curves piece of ch/nc_iw.tex is not included here.
\label{ch:pathdetail}

\section{More forms of the target, and three kinds of statement}\label{d:sec:kinds}

Section~\ref{route:sec:target} lists six forms (T1) to (T6) of RH, each with a parameter that measures how far from the line it is known. This section adds five more, each proved elsewhere in the three volumes or in the literature, and then sorts what is known in all of them into three kinds. The additions are the language of the Pólya frequency function, of the heat flow, of the mixing variable, and of the Hankel and Stieltjes continued fraction form of Chapter~\ref{ch:dbn}.

\begin{theorem}\label{d:thm:forms}
The following are equivalent to RH.
\begin{enumerate}
\item[(T7)] The Riemann density $F$, whose Fourier transform is $\xi(\frac12)/\xi(\frac12+v)$, is a Pólya frequency function.
\item[(T8)] In the de Bruijn and Newman family, $F_0$ is a Pólya frequency function; equivalently $\Lambda=0$.
\item[(T9)] The mixing variable $H$ with $\mathbb E\,e^{-uH}=\xi(\frac12)/\xi(\frac12+\sqrt{2u})$ exists and is a generalised gamma convolution.
\item[(T10)] The Hankel determinants $D_n$ and $D_n'$ of the Thorin ladder are positive for every $n$.
\item[(T11)] Every coefficient $c_n$ of the Stieltjes continued fraction of $\frac{\mathrm d}{\mathrm du}\log\xi(\frac12+\sqrt u)$ exists and is positive.
\end{enumerate}
\end{theorem}

\begin{proof}
Proposition~\ref{pol:prop:pf} and Theorem~\ref{p:thm:heat}, with the theorem of Rodgers and Tao \cite{RodgersTao} for the equality $\Lambda=0$ (Section~\ref{sec:ch21:flow}); Proposition~\ref{c:thm:wald}; Propositions~\ref{c:prop:grommer} and \ref{c:prop:string}, and Theorem~\ref{thm:ch12:inertia}, where the coefficients $c_n$ appear as the rates of a birth and death chain (Proposition~\ref{prop:ch21:verhankel} for the Hankel form).
\end{proof}

The Hankel conditions are positivity, not nonnegativity: nonnegative leading minors do not make a Hankel matrix positive semidefinite, and the criterion is positive definiteness of the centre Hankel matrices of Section~\ref{sec:ch21:hankel}.

Each route of Volumes I to III is placed against the statement about the zeros that it is equivalent to or implies in Table~\ref{d:tab:routemap}, with the quantifiers written out and the place in the volumes where the statement is proved or cited. The layout of the table is from an unrefereed note of Polson (November 2026). ``Exact'' means equivalent to a zero free statement by the results cited in the third column.

\begin{table}[t]
\centering\footnotesize
\begin{tabular}{@{}p{0.2\textwidth}p{0.52\textwidth}p{0.24\textwidth}@{}}
\toprule
Route & Relation to the zeros & Where \\
\midrule
Thorin ladder & Exact, level by level: for $\alpha\ge\frac12$, $G_\alpha$ is the Laplace transform of a GGC if and only if $\xi$ has no zero in $\re s>\alpha$, and then $U_\alpha$ is its Thorin measure. Unconditional for $\alpha\ge\frac78$, the theorem of \cite{OAI26a}; RH is $\alpha=\frac12$ & Theorem~\ref{bo:thm:level}, Corollary~\ref{bo:cor:level} \\
Cauchy deconvolution & Exact: given a zero free half plane $\re s>\theta$ and $\alpha\ge\theta$, RH holds if and only if $U_\alpha=P_{\alpha-1/2}*V$ for a locally finite signed $V$; then $V=U_{1/2}\ge0$ & Proposition~\ref{bo:prop:descent}, Theorem~\ref{tf:thm:deconv} \\
Chebyshev, prime side & The density at level $\alpha$ is written from the primes for $\alpha>\Theta=\sup\beta$, where the integration by parts converges; positivity on a single line is not shown to exclude zeros in $\re s>\alpha$ & Theorem~\ref{ct:thm:dictionary}, Remark~\ref{rem:ch2:oneline} \\
Hitting times & Exact: every rate of the chain of $\xi$ exists and is positive if and only if RH. A truncation to finitely many rates is a passage time whatever the zeros are, so a finite chain tests nothing until its rates are identified with those of $\xi$ & (T11) of Theorem~\ref{d:thm:forms}, Theorem~\ref{h:thm:crit} \\
Riemann density & Exact: $F$ is a Pólya frequency function if and only if RH. Each rung is a consequence of RH, not an equivalent; rungs one and two are certified & (T7) of Theorem~\ref{d:thm:forms}, Theorem~\ref{p:thm:m1} \\
Ferromagnetic, Lee and Yang & Exact: the Pólya law is ferromagnetic if and only if RH. An approximation route needs transforms that converge locally uniformly to that of $\xi$, or GGC reciprocal clocks that converge on $u\ge0$ to a limit continuous at $0$, and then each approximant carries its own zero statement & (T5) of Section~\ref{route:sec:target}, Chapter~\ref{chap:ferro}, \cite{Newman75} \\
de Bruijn Newman & Exact: RH holds if and only if $\Lambda\le0$, if and only if $\Lambda=0$; unconditionally $0\le\Lambda\le0.2$ \cite{RodgersTao,PT21} & (T8) of Theorem~\ref{d:thm:forms}, Corollary~\ref{cp:cor:known} \\
Hermite Biehler & Exact: for $\alpha>\frac12$, $E_\alpha(z)=\xi(\alpha-iz)$ is Hermite Biehler if and only if $\xi$ has no zero in $\re s\ge\alpha$; the strict inequality fails at $\alpha=\frac12$ even under RH & Theorem~\ref{mi:thm:paired} \\
Quasi RH ladders & A zero free half plane $\re s>\theta$ for $\zeta$ gives every rung $\alpha\ge\theta$; $\theta=\frac78$ is the theorem of \cite{OAI26a}. The theorem covers every Dirichlet $L$ function, $\zeta$ included & Corollary~\ref{bo:cor:level}, Section~\ref{qrh:sec:claims} \\
Weil positivity & Exact: $\sum_\rho\widehat h(t_\rho)\ge0$ for $h=g*\tilde g$ and every smooth $g$ of compact support if and only if RH \cite{Weil52}. Positivity for small support is a theorem \cite{Yoshida92}, and no zero free half plane is deduced from it here & (T4) of Section~\ref{route:sec:target}, Theorem~\ref{cm:thm:weil} \\
Li coefficients & Exact: $\lambda_n\ge0$ for every integer $n\ge1$ if and only if RH \cite{Li97,BL99}; a finite list of inequalities does not decide it & Section~\ref{ch:li} \\
B\'aez-Duarte & Exact: $\sum_{j=0}^k(-1)^j\binom kj/\zeta(2j+2)=O(k^{-3/4+\varepsilon})$ for every $\varepsilon>0$ if and only if RH \cite{BaezDuarte03} & Section~\ref{cw:sec:maslanka} \\
\bottomrule
\end{tabular}
\caption{Each route against the statement about the zeros. Routes marked exact are equivalent to RH or to the zero free half plane named in the row.}
\label{d:tab:routemap}
\end{table}


Each language turns RH into an infinite list of inequalities indexed by an integer: the rung $k$ of the ladder for $F$, the order $n$ of the Hankel determinants, the level $j$ of the chain. About any such list one can make three kinds of statement, and it clears the ground to say which kind has been achieved in each case.

\paragraph{Finite statements.} These assert the inequality for the indices up to some bound. Rungs one and two for $F$ are certified (Theorem~\ref{p:thm:m1}). The rates $c_1,\dots,c_{1299}$ of the chain of $\xi$ are positive, certified in ball arithmetic in the note and not repeated here (Section~\ref{sec:ch21:hankel}), so that the centre Hankel test passes to order $650$. Each finite statement is implied by a finite verification of zeros, and for rungs the detection rung theorem (Theorem~\ref{p:thm:detect}) shows that in the far field a few zeros off the line above that verification go unseen. Finite statements are therefore corollaries of the numerical verification of RH, obtained by a much harder route. Their value is that they test the machinery.

\paragraph{Asymptotic statements.} These assert the inequality for each fixed index in a limit or a range of some other parameter. Every rung of the ladder holds in the far field (Theorem~\ref{p:thm:farpfk}). Every rung holds for the heat flow family once $t\ge0.2$ (Theorem~\ref{p:thm:heat}). The rates of the smooth chain obey a law that makes them positive (Chapter~\ref{l:ch:bohr}, with the law \eqref{eq:ch21:ratelaw}, heuristic). The chain at level $\alpha$ is positive for every $\alpha>\frac78$ (Chapter~\ref{h:ch:chains}, \cite{OAI26a}). Asymptotic statements are genuine theorems about all indices at once, but each is made in a regime where the arithmetic is tame: far from the origin, after smoothing, with the fluctuation removed, or to the right of the strip.

\paragraph{Uniform statements.} A uniform statement asserts the inequality for all indices at the point that matters: all rungs for $F$ itself, all rates for the chain of $\xi$ itself. This is RH. No uniform statement is known in any of the languages, and we know of none in any equivalent formulation in the literature.

The gap between the second kind and the third is the whole problem. It has the same form in each language. There is a parameter that makes the problem easy when it is large, and RH sits at the end of its range:
\begin{center}\small
\begin{tabular}{@{}p{0.24\textwidth}p{0.3\textwidth}p{0.34\textwidth}@{}}
\toprule
parameter & easy regime & RH\\
\midrule
shift $L$ in a minor of $F$ & $L\to\infty$ for fixed $k$ & all $L$, for every $k$\\
heat time $t$ & $t\ge0.2$ & $t=0$\\
level $\alpha$ & $\alpha>\frac78$ & $\alpha=\frac12$\\
fluctuation $\varepsilon_j$ of the rates & $\varepsilon_j=0$, the gamma chain & the true $\varepsilon_j$, driven by the primes\\
\bottomrule
\end{tabular}
\end{center}
In each row the left end is a theorem and the right end is RH.

\section{How much room there is}\label{d:sec:room}

A natural hope is that RH holds with room to spare in some variable, so that a rough estimate would be enough. Three measurements give three different answers, and understanding why they differ is the best guide to what a proof must do.

\paragraph{In the heat time there is no room.} The heat flow gives a family in which the reality of the zeros is known on one side and false on the other.

\begin{proposition}\label{d:prop:closed}
RH holds if and only if $H_t$ has only real zeros for every $t>0$. For every $t<0$ the function $H_t$ has zeros off the real axis.
\end{proposition}

\begin{proof}
By the theorems of de Bruijn and Newman \cite{deBruijn,Newman1976} the set of $t$ for which $H_t$ has only real zeros is the closed half line $[\Lambda,\infty)$. So the condition for all $t>0$ is $\Lambda\le0$, which is RH. The second statement is $\Lambda\ge0$, the theorem of Rodgers and Tao (Section~\ref{sec:ch21:flow}).
\end{proof}

The first half says that RH is a limit of statements, each of which holds with some slack: for $t>0$ the zeros of $H_t$ are real and separated (known for $t\ge0.2$, and for $0<t<0.2$ conditional on RH). The second half removes the comfort: going an equal distance below $0$ destroys the property. The Riemann xi function lies exactly on the boundary of the set of functions with real zeros, along this family.
The consequence for proofs is concrete. An argument for RH that rests on a strict inequality which is stable under small perturbations of Riemann's kernel of the form $\Phi(v)\mapsto e^{tv^2/4}\Phi(v)$ would prove too much, since it would apply for small negative $t$. So whatever is proved about $\Phi$ must be an identity or an inequality that is sharp along the heat flow. The heat flow ladder shows the same thing in numbers: the first failing rung of $F_t$ is at most $4$ for $t=-4$, at most $135$ for $t=-0.25$ (Section~\ref{sec:ch21:hankel}), at most $4765$ for $t=-0.001$, and it tends to infinity as $t\uparrow0$ (computed in the note; the staircase of the Hankel test in Table~\ref{tab:ch21:nfail} is the same phenomenon in the language of Chapter~\ref{ch:dbn}).
% UNVERIFIED: 4 and 4765 from the draft (chC:52), not in W3; 135 is in Volume II (pair 135,136 collides at -0.2015); the pairs (135,136) and (4765,4766) match -gap^2/2 (U12 M1).

\paragraph{In the rates there appears to be a great deal of room.} Write the rates as $c_j=c_j^\Gamma(1+\varepsilon_j)$, with $c_j^\Gamma$ the rates of the gamma chain, approximated by \eqref{eq:ch21:ratelaw} with the constant $\frac74$ (heuristic). RH is the statement that $\varepsilon_j>-1$ for all $j$. In every computation $|\varepsilon_j|\le0.15$, to depth $2799$ (computed). A sequence that must stay above $-1$ and is observed within $0.15$ of zero seems to have room.
% UNVERIFIED: 0.15 and depth 2799 (chC:54); W3 V8 reached the gamma chain only to j=1399.

The control shows that the appearance is not to be trusted. Its fluctuation crosses $-1$ at four levels (the negative rates $c_{148}$, $c_{150}$, $c_{217}$ and $c_{219}$ of the Davenport and Heilbronn function \cite{DH36}, Section~\ref{sec:ch21:hankel}; computed). The typical size of the fluctuation says nothing about its extreme values. A zero off the line does not announce itself by a gradual growth of $|\varepsilon_j|$ over many levels. Its decisive effect, a negative rate, appears at the levels whose turning height is near the height of the zero.
% The 0.17 of the draft (chC:56) is dropped, see R_routes M5. The four levels agree with W3 V9 (148, 150, 217, 219). Locality of the response of the other rates is not claimed (U12 M6).

There is also a reason in principle for caution. The fluctuation is driven by the displacement of the zeros from their smooth positions, that is by the function $S(t)$, and $S(t)$ is unbounded. Its typical size grows like $\sqrt{\log\log t}$ \cite{Selberg46}, so slowly that no computation will ever see it grow. Whether the map from $S$ to $\varepsilon_j$ damps this growth is not known, and a bound $|\varepsilon_j|\le c<1$ for all $j$ would be comparable to RH itself.

\paragraph{In the density the room is exponentially small and located far out.} By the detection scale (Corollary~\ref{p:cor:scale}), a zero at distance $d$ from the line makes a minor of $F$ change sign with asymptotic spacing $\pi/d$ in the shift, so that the first negative minor is expected (heuristically) at shifts of that order, where the entries are of size $e^{-\gamma_1\pi/d}$ and smaller. The planted model of Section~\ref{e:sec:toy} shows this at small scale: with $d=0.4$ the second rung fails only from $x=8.75$, where the density is about $2\times10^{-6}$ (computed).% 8.75 and 2e-6 agree with ch/nc_v3_cosine.tex (table e:tab:toy).

\section{Two more requirements, and a warning}\label{d:sec:req}

The three filters of Section~\ref{route:sec:filters} are the first three requirements on any proof along this path: it must use the Euler product, it must see more than finitely many zeros (uniformity in the height), and it must not rest on a property that survives Gaussian damping, which by Proposition~\ref{d:prop:closed} means that it must be sharp along the heat flow. By the detection rung theorem a rung of index $k$ fails in the far field only for a zero off the line with fewer than $k$ real zeros below it; the corresponding statement for determinants is Proposition~\ref{prop:ch21:verhankel}, and none is proved for rates. So an argument that treats indices one at a time, or that needs more work for each new index, reproduces the numerical verification of RH and nothing more. The control has the gamma factor, a functional equation, a Dirichlet series with bounded coefficients, the same density of zeros and the same law for the smooth part of its chain, and it fails. This excludes arguments from the functional equation alone, from the growth and order of $\xi$ alone, and from the gamma factor alone, which is the lesson of Jackson's $q$ gamma function in Chapter~\ref{c:ch:cand}. Two further requirements follow.

\paragraph{4. It cannot be total positivity of the kernel.} Riemann's kernel is totally positive of order exactly four (computed, Proposition~\ref{c:thm:tp4}), and a kernel with an entire transform is a Pólya frequency function only if it is Gaussian (Proposition~\ref{c:prop:gauss}). The arithmetic does build positivity into the kernel, through the Euler factors, but that positivity stops at a finite order, and the order does not distinguish $\zeta$ from the control in the way required: the control and the $L$ function of the character modulo $4$ have the same order.

\paragraph{5. It cannot be approximation alone.} Functions with real zeros, and Pólya frequency functions, are closed under limits. But for any target there is always an approximating family that has the property: use the real parts of the zeros as knots (Proposition~\ref{c:prop:ord}). The content is in proving that a family with the property converges to $\xi$ itself and not to a neighbour. For the control the neighbour differs from the function by $7$ parts in $10^8$ in the density, and one of the two fails the ladder, at rung $44$ (computed).
% UNVERIFIED: 7e-8 and rung 44 (ch9:50).


\section{Closed roads}\label{d:sec:roads}

Two roads that look natural are closed by the results of the volumes, besides requirements 4 and 5. The first is finite computation, and is the first filter. The detection rung theorem (Theorem~\ref{p:thm:detect}) says that a zero off the line with $m$ simple real zeros below it makes $F$ fail rung $m+1$, while every configuration of at most $m$ points still passes in the far field. Since RH is verified to height $3\times10^{12}$ \cite{PT21}, every configuration of at most about $10^{13}$ points passes in the far field, and a rung beyond that cannot be reached. A failure would also be invisible in practice: by the detection scale (Corollary~\ref{p:cor:scale}) the failing minor is expected to turn negative at a shift of order $(\beta-\frac12)^{-1}$, where the entries of the matrix are of size $e^{-\pi\gamma_1/(\beta-1/2)}$.
The second is complete monotonicity alone. That $F$ is a scale mixture of normals, so that $H$ exists, is a single positivity condition. The control appears to satisfy it: its mixing density is positive to thirty orders of magnitude below its peak, and its moment ladder holds to $n=50$ (computed). Existence of $H$ is not the content of RH. That $H$ is a generalised gamma convolution is.
% UNVERIFIED: thirty orders, n=50 (ch9:52).

\section{Route A in detail: descent in the level}\label{d:sec:A}

The first route starts where the arithmetic is strong. For $\alpha>1$ the Thorin density at level $\alpha$ can be written from the primes, and its positivity comes from the positivity of the coefficients of $-\zeta'/\zeta$, that is of the von Mangoldt function. At $\alpha=1$ positivity is equivalent to the absence of zeros on the line $\re s=1$, the theorem of Hadamard and de la Vallée Poussin, whose proof uses the Euler product through the inequality $3+4\cos\theta+\cos2\theta\ge0$. Below $1$ the levels are tied together by the Cauchy semigroup: the density at level $\alpha$ is the Poisson smoothing of the density at any lower level, and going down is a deconvolution. This is (T1) to (T3) of Section~\ref{route:sec:target}.

Chapter~\ref{h:ch:chains} recasts this as a ladder of chains. The chain at level $\alpha$ has positive rates when every zero has real part below $\alpha$, and a regular chain with positive rates excludes zeros to the right of $\alpha$ except a pair at one ordinate placed symmetrically about $\alpha$ (Theorem~\ref{h:thm:crit}), and the rates at level $\alpha$ converge to the rates of the chain of the zeros as $\alpha\downarrow\frac12$ (Theorem~\ref{h:thm:limit}).
\paragraph{What is known.} Classical zero free regions are of the form $\re s>1-c/\log|t|$ and its refinements. They are regions that narrow with the height, and they do not contain any half plane $\re s>\alpha$ with $\alpha<1$. In the language of the ladder, level one is the top rung of the classical theory, and no lower rung follows from those regions. The descent from level one to $\frac78$ is done by the theorem of \cite{OAI26a} (see below); the open part starts below $\frac78$.

\paragraph{What the ladder adds.} Three things. It replaces a pointwise positivity, which can miss a zero near the line (Proposition~\ref{h:prop:dipole}), by positivity of the rates, which cannot. It gives a finite, certified test at each level, with a count: the number of negative rates is at most $4k$ for $k$ zeros to the right (Theorem~\ref{h:thm:upper}), and exactly governed by that number under the hypotheses of Theorem~\ref{h:thm:inertia}. And it shows a stability that was not expected: the rates of $\xi$ move by at most ten per cent between $\alpha=1$ and $\alpha=0.51$ (computed).
% UNVERIFIED: ten per cent (chC:90; ch8).

\paragraph{What Bondesson's letter adds.} Chapter~\ref{bo:ch:bondesson} removes the gamma factor from the descent exactly: the archimedean part of $U_\alpha$ moves between levels by the Poisson semigroup and is known in closed form, so only the arithmetic remainder, of mean zero, has to be carried down. This improves linear deconvolution from the levels $\frac34$ and $\frac78$ by a factor of $9$ to $62$ on exact data and of $1.4$ to $4.9$ on the data of the speed study, and Proposition~\ref{bo:prop:descent} states exactly what descent from a zero free level must show. Below level $1$ it does not itself supply positivity, since the arithmetic part stops being infinitely divisible there; positivity down to $\frac78$ comes from the theorem of \cite{OAI26a}.

\paragraph{A criterion at one point.} Section~\ref{bo:sec:ledger} turns the descent into a single inequality: there are no zeros to the right of $\alpha$ exactly when a Fejér weighted sum over prime powers up to $Y$ stays below three closed form terms for all large $Y$ (Theorem~\ref{bo:thm:fejer}). On the critical line the slack is, under RH, a $2\sin^2$ sum over the zeros of size $1/\log Y$, while each side is of size $\sqrt Y/\log Y$. This is route A in its most compact form, and it shows the price exactly. The inequality is a theorem at level $1$, for every cutoff (Corollary~\ref{bo:cor:mertens}); above the line it has a margin linear in the distance; on the line it has none, since the slack returns arbitrarily close to zero (Proposition~\ref{bo:prop:margin}). So route A, too, ends at a statement with no room, as requirement 3 said it must. Theorem~\ref{bo:thm:osc} also joins route A to route B: the slack is the cosine transform of the spectral measure $\nu$ of the chain, whose moments give the rates.

\paragraph{The quasi Riemann half plane.} The zero free half plane $\re s>\frac78$ is the theorem of \cite{OAI26a} (Section~\ref{qrh:sec:claims}); the $\frac{11}{12}$ manuscript \cite{OAI26b} is separate and is not in the Lean catalogue recorded there. The theorem is precisely a first step of route A: every instrument of this volume is valid at every level above $\frac78$.
\paragraph{What a step would be.} A step is a proof that the chain at some level $\alpha_0\le\frac78$ is positive. Its content is a zero free half plane beyond the theorem of \cite{OAI26a}, which would be a major theorem in its own right. The realistic first targets are below that. One is a proof of the known rung in the language of the rates: derive the positivity of the chain at $\alpha=1$ directly from the Euler product. This shows the form in which arithmetic positivity enters a continued fraction; Proposition~\ref{qt:prop:levelone} gives such a derivation, using only that $\zeta$ has no zeros with $\Re s>1$ and no information about the nontrivial zeros. The second is to explain the stability: a uniform bound on the logarithm of a rate, from a level $\alpha>\frac78$, would carry positivity down the ladder, but it is needed for all $j$ at once (requirement 2).


\section{Route C in detail: limits and the heat flow}\label{d:sec:C}

The third route approaches $\xi$ from the side on which the answer is known. For $t\ge0.2$ every rung of the ladder holds for $F_t$, and by Proposition~\ref{d:prop:closed} RH is the statement that this persists for every $t>0$.

\paragraph{What is known.} The upper bound $\Lambda\le0.2$ \cite{PT21} is obtained by combining a numerical verification of RH up to a height with an analytic argument above that height, where the Gaussian factor makes the zeros of $H_t$ regular. The height needed grows rapidly as $t$ decreases, and each improvement of the bound has required a larger verification. In the other direction, each pair of close zeros of $\zeta$ gives a negative lower bound for $\Lambda$ by the theorem of Csordas, Smith and Varga \cite{CSV94}, and the collision time of Lehmer's pair $\gamma_{6709},\gamma_{6710}$, itself a lower bound for $\Lambda$, is $-7.11\times10^{-4}$ (Section~\ref{sec:ch21:collisions}).% -7.11e-4 and the pair (6709, 6710) are in ch/ch21.tex (collisions). The draft says the first 8900 zeros; the zeros known to the draft number 6703 in chD, so the range is left out.

\paragraph{What the volume adds.} The heat flow ladder attaches a rung to each $t<0$: the function $F_t$ fails by rung $m(t)+1$ when the lowest nonreal layer is one simple quadruple of zeros, where $m(t)$ counts the real zeros of $H_t$ below it (Theorem~\ref{p:thm:heat}). The sharpness tests at $t=-4$ and $t=-2$ confirm that the failure occurs at the predicted rung and with the predicted period in the shift (computed).
% UNVERIFIED: the sharpness tests at t=-4 and t=-2 (chC:130) are not in the verification files. So the family $F_t$ is a set of controls of adjustable difficulty, all built from $\xi$ itself and all carrying its arithmetic. This is useful in a way that the Davenport and Heilbronn function is not: a mechanism can be tested for sharpness along the flow, which is requirement 3. In Chapter~\ref{ch:dbn} the Hankel test along the flow already locates, for $t<0$, the index $n_{\rm fail}(t)$ at which it first fails (Table~\ref{tab:ch21:nfail}), and Theorem~\ref{thm:ch13:perturb} says why a perturbation argument cannot decide the case $t=0$.

\paragraph{What a step would be.} A proof that $F_t$ passes every rung up to $m(t)$ for $t<0$ would make the failing rung exact and would be a first uniform statement, in $t$, about the ladder. On the side of weak limits, the task is to find a family of laws that is known to have the property, hitting times or ferromagnetic laws, and that converges to the law of Pólya's variable. Requirement 5 says that convergence is the entire content.
